\documentclass[a4paper,11pt]{article}

\usepackage[T1]{fontenc}
\usepackage[margin=2.5cm]{geometry}
\usepackage{booktabs}
\usepackage{listings}
\usepackage{xcolor}
\usepackage{hyperref}

\lstset{
  basicstyle=\ttfamily\small,
  backgroundcolor=\color{gray!10},
  frame=single,
  framerule=0pt,
  xleftmargin=6pt,
  xrightmargin=6pt,
  breaklines=true,
}

\title{\textsf{beamerthemeInsiDark} --- Minimal NU dark Beamer theme\\[4pt]
  \large Version 1.0}
\author{Pietro Brach del Prever\\
  \small\href{mailto:brachdelprever.p@northeastern.edu}{brachdelprever.p@northeastern.edu}}
\date{2026/05/07}

\begin{document}
\maketitle

\begin{abstract}
	\textsf{beamerthemeInsiDark} is a minimal dark Beamer theme for
	Northeastern University / IINS presentations.  It features a midnight-navy
	header strip with a textured background image, the ColorBrewer
	\texttt{Paired} palette for content, Fira Sans typography, and convenience
	commands for section separators, closing slides, and event metadata.
\end{abstract}

\section{Requirements}

\begin{itemize}
	\item \textsf{FiraSans} and \textsf{FiraMono} ---
	      \texttt{tlmgr install fira}
	\item \textsf{tikz}, \textsf{pgfplots} (with the \texttt{colorbrewer}
	      library), \textsf{graphicx} (part of standard TeX Live / MacTeX)
	\item \textsf{currfile} (part of standard TeX Live / MacTeX)
	\item \texttt{assets/bg-dark.png} and \texttt{assets/logo-insi.png} ---
	      bundled with this package
\end{itemize}

The theme finds its own \texttt{assets/} directory automatically via
\textsf{currfile}, so no assets need to be copied into the document project.

\section{Installation}

Symlink the package directory into the personal TEXMF tree (macOS/MacTeX):

\begin{lstlisting}
ln -s /path/to/beamer-template \
      ~/Library/texmf/tex/latex/beamerthemeInsiDark
\end{lstlisting}

After that, \verb|\usetheme{InsiDark}| resolves from any document on the
system.  Verify with:

\begin{lstlisting}
kpsewhich beamerthemeInsiDark.sty
\end{lstlisting}

\section{Loading the theme}

\begin{lstlisting}
\documentclass[aspectratio=169, 10pt]{beamer}
\usetheme{InsiDark}
\end{lstlisting}

\section{Metadata commands}

Standard Beamer metadata (\verb|\title|, \verb|\subtitle|, \verb|\author|,
\verb|\institute|, \verb|\date|) is used as usual.  The theme adds one extra
command:

\begin{quote}
	\verb|\event{|\textit{name}\verb|}|
\end{quote}

Sets the conference or event name shown on the title page below the institute
line.

\section{Slide commands}

\subsection{Title slide}

The title page must be placed inside a group that sets the full-bleed dark
background canvas:

\begin{lstlisting}
{
  \setbeamertemplate{background}{}
  \setbeamertemplate{background canvas}{\InsiDarkBg}
  \setbeamertemplate{footline}{}
  \begin{frame}[plain]
    \titlepage
  \end{frame}
}
\end{lstlisting}

\subsection{\texttt{\textbackslash sectionframe}}

Inserts a dark full-bleed slide that displays the current section name
centered in white.  Call it immediately after \verb|\section{...}|:

\begin{lstlisting}
\section{Results}
\sectionframe
\end{lstlisting}

\subsection{\texttt{\textbackslash closingframe\{msg\}\{contact\}}}

Inserts a dark closing slide with a large message and a contact line:

\begin{lstlisting}
\closingframe{Thank you}{your@email.edu}
\end{lstlisting}

\subsection{\texttt{\textbackslash InsiDarkBg}}

An \verb|\includegraphics| call that tiles the full-bleed background image
to the slide dimensions.  Use it as the argument to
\verb|\setbeamertemplate{background canvas}{...}| for any custom dark slide.

\subsection{\texttt{\textbackslash nuLogos}}

Overlays the INSI monogram logo at the top-right corner of the current slide
using a TikZ \texttt{remember picture, overlay} node.  Used automatically on
the title slide; can be called on any custom dark frame.

\section{Color palette}

Content colors come from the ColorBrewer \texttt{Paired} 12-color set, loaded
via the \textsf{pgfplots} \texttt{colorbrewer} library.  They are available as
\texttt{Paired-A} through \texttt{Paired-L} and can be used anywhere via
\verb|\textcolor{Paired-B}{...}|, \verb|\colorbox{Paired-D}{...}|, etc.
Light/dark pairs: A/B (blue), C/D (green), E/F (red), G/H (orange),
I/J (purple), K/L (yellow / brown).

The theme also defines a handful of reserved neutrals and a title-slide
accent:

\begin{center}
	\renewcommand{\arraystretch}{1.3}
	\begin{tabular}{lll}
		\toprule
		Name              & Hex               & Role                                           \\
		\midrule
		\texttt{InsiDark} & \texttt{\#01091E} & Title / section / header background            \\
		\texttt{themeMid} & \texttt{\#666666} & Footer / caption text                          \\
		\texttt{themeBg}  & \texttt{\#F4F4F4} & Neutral block body background                  \\
		\texttt{nuGold}   & \texttt{\#FDBB45} & Title-slide subtitle accent (title slide only) \\
		\bottomrule
	\end{tabular}
\end{center}

Default block colors are \texttt{Paired-B} (info / standard block),
\texttt{Paired-F} (alert), and \texttt{Paired-D} (example).  Itemize and
enumerate bullets are plain black.

\section{Full example}

The file \texttt{main.tex} (included in this package) is a complete working
template demonstrating all slide types.  It compiles with:

\begin{lstlisting}
pdflatex main.tex
\end{lstlisting}

The minimal skeleton:

\begin{lstlisting}
\documentclass[aspectratio=169, 10pt]{beamer}
\usepackage{microtype}
\usetheme{InsiDark}

\title{Presentation Title}
\subtitle{A Concise Subtitle}
\author[P.\ Brach del Prever]{Pietro Brach del Prever}
\institute[NU / IINS]{Northeastern University \\
                      Institute for Intelligent Networked Systems}
\event{IEEE INFOCOM 2026 Workshop}
\date{May 2026}

\begin{document}

% Title slide
{
  \setbeamertemplate{background}{}
  \setbeamertemplate{background canvas}{\InsiDarkBg}
  \setbeamertemplate{footline}{}
  \begin{frame}[plain]\titlepage\end{frame}
}

\section{Introduction}
\sectionframe

\begin{frame}{Motivation}
  \begin{itemize}
    \item First key point.
    \item Second key point.
  \end{itemize}
\end{frame}

\begin{frame}{Blocks}
  \begin{block}{Definition}
    Formal statement.
  \end{block}
  \begin{alertblock}{Key Result}
    Main contribution.
  \end{alertblock}
  \begin{exampleblock}{Example}
    Concrete illustration.
  \end{exampleblock}
  \begin{block}{}
    An empty title drops the title bar: a body-only callout.
  \end{block}
  \begin{narrow}[0.5\textwidth]
    \begin{block}{Narrow}
      \texttt{narrow} tightens anything sized from \texttt{\textbackslash textwidth}.
    \end{block}
  \end{narrow}
\end{frame}

\closingframe{Thank you}{your@email.edu}

\end{document}
\end{lstlisting}

\section{License}

This work is released under the \LaTeX\ Project Public License,
version~1.3c or later.
See \url{https://www.latex-project.org/lppl.txt}.

\end{document}
